\section{The guarded translation and its intended models}
\label{sec:translation}

This section defines the guarded translation of \Frag\ into untyped FOL,
the axiom set it emits, and --- the heart of the soundness argument --- a
family of untyped models validating the emitted axioms.

\subsection{Target signature and translation}

\begin{definition}[Target signature]
  \label{def:target-sig}
  The untyped signature $\Sig$ has:
  \begin{itemize}
  \item a function symbol $\hat h$ of arity $k_h + m_h$ for every
    $h \in \Constr \cup \Fun$;
  \item a function symbol $\hat I$ of arity $k_I$ for every $I \in \Ind$;
  \item a constant $\TypeC$;
  \item a predicate symbol $\hat p$ of arity $k_p + m_p$ for every
    $p \in \Pred$, and a binary predicate symbol $\T$.
  \end{itemize}
  We reuse \Frag's type variables and term variables as FOL variables.
\end{definition}

\begin{definition}[Translation]
  \label{def:translation}
  The term translation $\enc{\cdot}$ maps \Frag-terms and \Frag-types to
  $\Sig$-terms:
  \[
    \enc{x} \defeq x, \qquad
    \enc{\alpha} \defeq \alpha, \qquad
    \enc{I(\tup\sigma)} \defeq \hat I(\enc{\tup\sigma}), \qquad
    \enc{h\langle\tup\sigma\rangle(\tup t)} \defeq
      \hat h(\enc{\tup\sigma}, \enc{\tup t}).
  \]
  The formula translation is defined by
  \begin{align*}
    \enc{p\langle\tup\sigma\rangle(\tup t)}
      &\defeq \hat p(\enc{\tup\sigma}, \enc{\tup t})
    & \enc{t \eqx{\tau} s} &\defeq \enc{t} = \enc{s} \\
    \enc{\forall x{:}\tau.\ \phi}
      &\defeq \forall x.\ \T(x, \enc{\tau}) \imp \enc{\phi}
    & \enc{\exists x{:}\tau.\ \phi}
      &\defeq \exists x.\ \T(x, \enc{\tau}) \wedge \enc{\phi} \\
    \enc{\forall \alpha.\ \phi}
      &\defeq \forall \alpha.\ \T(\alpha, \TypeC) \imp \enc{\phi}
  \end{align*}
  and homomorphically on $\bot$, $\top$, $\wedge$, $\vee$, $\imp$.
\end{definition}

\begin{definition}[Emitted axioms]
  \label{def:axioms}
  The \emph{translation of the environment} is
  $\Ax \defeq \mathrm{TA} \cup \mathrm{ID} \cup
  \enc{\mathrm{Prem}(\Env)} \cup \enc{\StructAx_{\exists}}$, where:
  \begin{enumerate}[label=(T\arabic*)]
  \item \emph{Typing axioms} $\mathrm{TA}$: for every
    $h \in \Constr \cup \Fun$,
    \[
      \forall \tup\alpha\,\tup x.\
      \bigwedge_{j} \T(\alpha_j, \TypeC) \imp
      \bigwedge_{i} \T(x_i, \enc{\tau^h_i}) \imp
      \T\bigl(\hat h(\tup\alpha,\tup x),\ \enc{\tau^h_0}\bigr),
    \]
    (iterated implication, associated to the right) and for every
    $I \in \Ind$:
    $\forall\tup\alpha.\ \bigwedge_j \T(\alpha_j,\TypeC) \imp
    \T(\hat I(\tup\alpha), \TypeC)$.
  \item \emph{Unguarded injectivity and discrimination} $\mathrm{ID}$: for
    every constructor $c$ of $I$,
    \[
      \forall \tup\alpha\,\tup x\,\tup\alpha'\,\tup x'.\
      \hat c(\tup\alpha,\tup x) = \hat c(\tup\alpha',\tup x') \imp
      \textstyle\bigwedge_j \alpha_j = \alpha'_j \wedge
      \bigwedge_i x_i = x'_i,
    \]
    and for every pair $c \neq c'$ of constructors of the same $I$:
    $\forall \tup\alpha\,\tup x\,\tup\alpha'\,\tup x'.\
    \hat c(\tup\alpha,\tup x) \neq \hat c'(\tup\alpha',\tup x')$.
    These are deliberately \emph{unguarded} and quantify the two sides'
    type arguments independently, as hammer translations emit them.
  \item \emph{Translated premises and structural sentences}: the
    $\enc{\cdot}$-images of all sentences in $\mathrm{Prem}(\Env)$ and of
    the exhaustiveness sentences (S3) and predicate
    introduction/inversion sentences of $\StructAx$ (written
    $\StructAx_\exists$; the guarded images of (S1)--(S2) are logical
    consequences of $\mathrm{ID}$ and need not be emitted).
  \end{enumerate}
  Separately, the \emph{typehood-completion axioms} are
  $\Comp \defeq \mathrm{Cp} \cup \mathrm{Ct}$ with
  \begin{enumerate}[label=(C\arabic*)]
  \item $\mathrm{Cp}$: for every $p \in \Pred$, every $j \le k_p$ and every
    $i \le m_p$,
    \[
      \forall\tup\alpha\,\tup x.\
      \hat p(\tup\alpha,\tup x) \imp \T(\alpha_j, \TypeC),
      \qquad
      \forall\tup\alpha\,\tup x.\
      \hat p(\tup\alpha,\tup x) \imp \T(x_i, \enc{\tau^p_i});
    \]
  \item $\mathrm{Ct}$: for every $I \in \Ind$ and $j \le k_I$:
    $\forall\tup\alpha\,z.\ \T(z, \hat I(\tup\alpha)) \imp
    \T(\alpha_j, \TypeC)$.
  \end{enumerate}
  Given a goal sentence $\phi_0$, the \emph{translated problem} is
  $\Ax \cup \Comp \cup \{\neg\enc{\phi_0}\}$.
\end{definition}

\subsection{The intended models \texorpdfstring{$\MN$}{M(N)}}
\label{subsec:intended}

Fix a ground structure $\NN \models \Th$.  We build an untyped
$\Sig$-structure whose domain contains, besides denotations for all ground
types and all their inhabitants, explicit \emph{junk} elements representing
ill-typed applications.  Junk is generated \emph{freely}, which is what
validates the unguarded axioms $\mathrm{ID}$.

\begin{definition}[Domain]
  \label{def:domain}
  Let
  \[
    \Dom_0 \defeq
    \underbrace{\GTy}_{\text{type elements}} \;\uplus\;
    \underbrace{\{\ast\}}_{\text{the element } \TypeC} \;\uplus\;
    \underbrace{\{(\tau, a) \mid \tau \in \GTy,\ a \in \NN_\tau\}}_{
      \text{tagged carrier elements}} .
  \]
  For a function symbol $g \in \Sig$ of arity $n$ and a tuple
  $\tup d \in (\Dom_0 \uplus \Junk)^n$ define the \emph{sort check}
  $\mathrm{ok}_g(\tup d)$:
  \begin{itemize}
  \item $\mathrm{ok}_{\hat I}(\tup d)$ iff $d_j \in \GTy$ for all
    $j \le k_I$;
  \item $\mathrm{ok}_{\hat h}(\tup d)$ (for $h \in \Constr \cup \Fun$) iff
    $d_j \in \GTy$ for $j \le k_h$ --- write $\tau_j \defeq d_j$ --- and
    $d_{k_h+i} = \bigl(\tau^h_i[\tup\tau/\tup\alpha],\, a_i\bigr)$ for some
    $a_i \in \NN_{\tau^h_i[\tup\tau/\tup\alpha]}$, for all $i \le m_h$;
  \item $\mathrm{ok}_{\TypeC}() $ always.
  \end{itemize}
  The junk universe $\Junk$ is the least set of formal pairs closed under:
  if $g \in \Sig$ is a function symbol of arity $n \ge 1$,
  $\tup d \in (\Dom_0 \uplus \Junk)^n$ and $\neg\mathrm{ok}_g(\tup d)$,
  then $(g, \tup d) \in \Junk$.  (Formally,
  $\Junk = \bigcup_n \Junk_n$ with $\Junk_0 = \emptyset$ and
  $\Junk_{n+1}$ the set of such pairs over $\Dom_0 \uplus \Junk_n$; the
  union is closed because $\mathrm{ok}_g$ never holds when some argument
  lies in $\Junk$.)  The four classes of elements --- ground types, $\ast$,
  tagged pairs, junk pairs --- are pairwise disjoint (formally, tag them).
  Finally $\Dom \defeq \Dom_0 \uplus \Junk$; note $\Dom \neq \emptyset$
  since $\ast \in \Dom$.
\end{definition}

\begin{definition}[The structure $\MN$]
  \label{def:MN}
  $\MN$ is the $\Sig$-structure with domain $\Dom$ and:
  \begin{itemize}
  \item $\TypeC^{\MN} \defeq \ast$;
  \item $\hat I^{\MN}(\tup d) \defeq I(\tup d) \in \GTy$ if
    $\mathrm{ok}_{\hat I}(\tup d)$, and $(\hat I, \tup d) \in \Junk$
    otherwise;
  \item for $h \in \Constr \cup \Fun$:
    $\hat h^{\MN}(\tup d) \defeq
    \bigl(\tau^h_0[\tup\tau/\tup\alpha],\
    h_{\tup\tau}^{\NN}(\tup a)\bigr)$ if $\mathrm{ok}_{\hat h}(\tup d)$
    holds with $\tup\tau, \tup a$ as in \cref{def:domain}, and
    $(\hat h, \tup d) \in \Junk$ otherwise;
  \item $\T^{\MN}(d, u)$ iff either $u = \ast$ and $d \in \GTy$, or
    $u \in \GTy$ and $d = (u, a)$ for some $a \in \NN_u$;
  \item $\hat p^{\MN}(\tup d)$ iff $\mathrm{ok}$ holds in the evident
    sense --- $d_j \in \GTy$ for $j \le k_p$ (write $\tup\tau$) and
    $d_{k_p+i} = (\tau^p_i[\tup\tau/\tup\alpha], a_i)$ --- \emph{and}
    $p^{\NN}_{\tup\tau}(\tup a)$ holds.
  \end{itemize}
  All interpretations are total and well defined: the sort check determines
  $\tup\tau$ and $\tup a$ uniquely (tags are part of the elements), and the
  junk branch applies exactly when it fails.
\end{definition}

Note the two design decisions embodied in \cref{def:MN}.  First,
$\hat p^{\MN}$ is \emph{false-extended}: it fails on every argument tuple
violating $p$'s type discipline.  This is the semantic core of guard
omission --- it is the model-theoretic counterpart of ``an ill-typed
application is not a provable proposition'' --- and it makes the completion
axioms true (\cref{thm:base-soundness}).  Second, junk is \emph{free}: an
ill-typed application denotes the formal tree recording its head and
arguments, so distinct ill-typed applications denote distinct elements.
This validates the unguarded injectivity and discrimination axioms, giving
a precise sense to the folklore justification that they are ``true in the
term model''.

\subsection{Denotation and relativization}

Throughout, let $\phi, t, \tau$ be well formed in a context $(V;\Gamma)$,
let $\theta : V \arr \GTy$ be a ground type instantiation, and let $v$ be a
\emph{proper valuation}: a sort-respecting assignment
$v(x) \in \NN_{\Gamma(x)\theta}$ for term variables.  Define the FOL
valuation $w_{\theta,v}$ on $\MN$ by
$w_{\theta,v}(\alpha) \defeq \theta(\alpha)$ for $\alpha \in V$ and
$w_{\theta,v}(x) \defeq (\Gamma(x)\theta,\, v(x))$.

\begin{lemma}[Denotation]
  \label{lem:denotation}
  For every type $\sigma \in \Ty{V}$:
  $\den{\enc{\sigma}}{\MN, w_{\theta,v}} = \sigma\theta \in \GTy$.
  For every term $t$ with $\wfj{V}{\Gamma}{t : \tau}$:
  \[
    \den{\enc{t}}{\MN, w_{\theta,v}} =
    \bigl(\tau\theta,\ \den{t\theta}{\NN, v}\bigr),
  \]
  where $\den{t\theta}{\NN,v}$ is the many-sorted denotation of the ground
  instance $t\theta$ (\cref{def:ms-structures}).
\end{lemma}

\begin{proof}
  Both claims by induction on the syntax.  Types: a variable $\alpha$
  denotes $\theta(\alpha)$; for $I(\tup\sigma)$ the induction hypothesis
  gives arguments in $\GTy$, so $\mathrm{ok}_{\hat I}$ holds and the value
  is $I(\tup\sigma\theta) = I(\tup\sigma)\theta$.  Terms: a variable $x$
  denotes $w_{\theta,v}(x) = (\Gamma(x)\theta, v(x))$ and
  $\tau = \Gamma(x)$ by the typing rule.  For
  $t = h\langle\tup\sigma\rangle(\tup t)$ with
  $\wfj{V}{\Gamma}{t_i : \tau^h_i[\tup\sigma/\tup\alpha]}$: by the
  induction hypotheses the type arguments denote
  $\tup\sigma\theta \in \GTy^{k_h}$ and the term arguments denote
  $\bigl(\tau^h_i[\tup\sigma/\tup\alpha]\theta,\,
  \den{t_i\theta}{\NN,v}\bigr) =
  \bigl(\tau^h_i[\tup\sigma\theta/\tup\alpha],\,
  \den{t_i\theta}{\NN,v}\bigr)$ (substitution composition, as in
  \cref{lem:instantiation}).  Hence $\mathrm{ok}_{\hat h}$ holds with
  $\tup\tau = \tup\sigma\theta$, and the value is
  $\bigl(\tau^h_0[\tup\sigma\theta/\tup\alpha],\,
  h^{\NN}_{\tup\sigma\theta}(\dots)\bigr)$, which is exactly
  $(\tau\theta, \den{t\theta}{\NN,v})$ by the definition of many-sorted
  denotation of the instance.
\end{proof}

\begin{lemma}[Relativization]
  \label{lem:relativization}
  For every type-quantifier-free formula $\phi$ well formed in
  $(V;\Gamma)$, every $\theta$ and proper $v$ as above:
  \[
    \MN, w_{\theta,v} \models \enc{\phi}
    \quad\text{iff}\quad
    \NN, v \models \phi\theta .
  \]
  Consequently, for every sentence $\phi_0$ (with prenex type
  quantifiers): $\MN \models \enc{\phi_0}$ iff $\NN \entg \phi_0$.
\end{lemma}

\begin{proof}
  Induction on $\phi$.

  \emph{Atoms.}  $\enc{p\langle\tup\sigma\rangle(\tup t)} =
  \hat p(\enc{\tup\sigma},\enc{\tup t})$.  By \cref{lem:denotation} the
  arguments denote $\tup\sigma\theta \in \GTy$ and the tagged pairs
  $(\tau^p_i[\tup\sigma\theta/\tup\alpha], \den{t_i\theta}{\NN,v})$, so
  $\mathrm{ok}$ holds and, by \cref{def:MN}, the atom is true in $\MN$ iff
  $p^{\NN}_{\tup\sigma\theta}(\dots)$ holds, i.e.\ iff
  $\NN,v \models (p\langle\tup\sigma\rangle(\tup t))\theta$.

  \emph{Equations.}  $\enc{t} = \enc{s}$ holds iff the two tagged pairs are
  equal; since both carry the same tag $\tau\theta$
  (\cref{lem:denotation}, using that $t$ and $s$ have the same type
  $\tau$), this is equivalent to
  $\den{t\theta}{\NN,v} = \den{s\theta}{\NN,v}$, i.e.\ to
  $\NN,v \models (t \eqx{\tau} s)\theta$.

  \emph{Connectives.}  Immediate from the induction hypotheses, as both
  semantics interpret them classically.

  \emph{Term quantifiers.}  Consider
  $\enc{\forall x{:}\tau.\ \phi} = \forall x.\ \T(x,\enc{\tau}) \imp
  \enc{\phi}$.  By \cref{lem:denotation},
  $\den{\enc{\tau}}{\MN,w_{\theta,v}} = \tau\theta \in \GTy$, so for
  $d \in \Dom$ the guard $\T(d, \tau\theta)$ holds iff $d = (\tau\theta, a)$
  with $a \in \NN_{\tau\theta}$.  Hence the $\MN$-side quantification,
  restricted by the guard, ranges exactly over
  $\{(\tau\theta,a) \mid a \in \NN_{\tau\theta}\}$; for each such $d$ the
  extended valuation is $w_{\theta, v[x \mapsto a]}$ and the induction
  hypothesis applies.  This matches the $\NN$-side clause ``for all
  $a \in \NN_{\tau\theta}$'', including the empty-carrier case (both sides
  vacuously true).  The existential case is dual (both sides false over an
  empty carrier).

  \emph{Sentences.}  For
  $\phi_0 = \forall\seq{\alpha}{n}.\psi$:
  $\enc{\phi_0} = \forall\tup\alpha.\ \T(\alpha_1,\TypeC) \imp \dots$, and
  $\T^{\MN}(d, \ast)$ holds iff $d \in \GTy$.  So
  $\MN \models \enc{\phi_0}$ iff $\MN, [\tup\alpha \mapsto \tup\tau]
  \models \enc{\psi}$ for all $\tup\tau \in \GTy^n$, which by the main
  claim (with $\theta = [\tup\tau/\tup\alpha]$, empty $\Gamma$-part) is
  equivalent to $\NN \models \psi[\tup\tau/\tup\alpha]$ for all
  $\tup\tau$, i.e.\ to $\NN \entg \phi_0$.
\end{proof}

\subsection{Base soundness}

\begin{theorem}[Model correctness]
  \label{thm:base-soundness}
  For every ground structure $\NN \models \Th$:
  \ $\MN \models \Ax \cup \Comp$.
\end{theorem}

\begin{proof}
  We check each axiom class against \cref{def:MN}; fix an arbitrary
  valuation $w$ into $\Dom$.

  \emph{(T1) Typing axioms.}  Assume the guards: $w(\alpha_j) \in \GTy$
  for all $j$ (so they determine $\theta$) and
  $\T(w(x_i), \den{\enc{\tau^h_i}}{w})$ for all $i$.  Since
  $w(\tup\alpha) = \tup\tau \in \GTy$, the first part of
  \cref{lem:denotation} gives
  $\den{\enc{\tau^h_i}}{w} = \tau^h_i[\tup\tau/\tup\alpha]$ (the lemma's
  proof uses only the values of $w$ on $\ftv{\tau^h_i}$, all in $\GTy$).
  So each $w(x_i)$ is a tagged pair
  $(\tau^h_i[\tup\tau/\tup\alpha], a_i)$; thus $\mathrm{ok}_{\hat h}$
  holds, the value of $\hat h(\tup\alpha, \tup x)$ under $w$ is
  $(\tau^h_0[\tup\tau/\tup\alpha], h^{\NN}_{\tup\tau}(\tup a))$, and the
  conclusion $\T$-atom holds.  The $\hat I$ axioms are analogous
  ($I(\tup\tau) \in \GTy$ and $\T(I(\tup\tau), \ast)$ holds).

  \emph{(T2) Injectivity.}  Let
  $d \defeq \hat c^{\MN}(w(\tup\alpha), w(\tup x))$ and
  $d' \defeq \hat c^{\MN}(w(\tup\alpha'), w(\tup x'))$ and assume $d = d'$.
  Three cases.  (i) Both sort checks succeed: then
  $d = (I(\tup\tau), c^{\NN}_{\tup\tau}(\tup a))$ and
  $d' = (I(\tup\tau'), c^{\NN}_{\tup\tau'}(\tup a'))$; equality of the
  tags gives $I(\tup\tau) = I(\tup\tau')$ in $\GTy$, whence
  $\tup\tau = \tup\tau'$ because ground types are syntax (a first-order
  term algebra, in which constructors are injective); then
  $c^{\NN}_{\tup\tau}(\tup a) = c^{\NN}_{\tup\tau}(\tup a')$, and
  $\NN \entg$ (S1) at instance $\tup\tau$ yields $\tup a = \tup a'$; hence
  $w(\alpha_j) = \tau_j = \tau'_j = w(\alpha'_j)$ and
  $w(x_i) = (\tau^c_i[\tup\tau/\tup\alpha], a_i) = w(x'_i)$.  (ii) Both
  fail: $d = (\hat c, w(\tup\alpha)w(\tup x))$ and
  $d' = (\hat c, w(\tup\alpha')w(\tup x'))$ are formal pairs; their
  equality gives componentwise equality of the argument tuples, which is
  the conclusion.  (iii) Exactly one fails: then $d$ is a tagged pair and
  $d'$ a junk pair (or vice versa), contradicting $d = d'$; the
  implication holds vacuously.

  \emph{(T2) Discrimination.}  With $d, d'$ as above for $c \neq c'$ of
  the same $I$: if both sort checks succeed and the tags agree, then
  $\tup\tau = \tup\tau'$ as before and
  $\NN \entg$ (S2) at $\tup\tau$ gives
  $c^{\NN}_{\tup\tau}(\tup a) \neq c'^{\NN}_{\tup\tau}(\tup a')$; if the
  tags differ, $d \neq d'$ outright.  If both fail, the junk pairs have
  distinct heads $\hat c \neq \hat c'$.  Mixed cases are distinct by
  disjointness.

  \emph{(T3) Translated sentences.}  Each is $\enc{\phi}$ for a sentence
  $\phi$ with $\NN \entg \phi$ (premises and the retained structural
  sentences are in $\Th$); apply \cref{lem:relativization}.

  \emph{(C1) Predicate completion.}  If
  $\hat p^{\MN}(w(\tup\alpha), w(\tup x))$ holds then, by
  \cref{def:MN}, $w(\alpha_j) \in \GTy$ --- so
  $\T(w(\alpha_j), \ast)$ holds --- and
  $w(x_i) = (\tau^p_i[\tup\tau/\tup\alpha], a_i)$; since
  $\den{\enc{\tau^p_i}}{w} = \tau^p_i[\tup\tau/\tup\alpha]$
  (\cref{lem:denotation}, as in (T1)), the atom
  $\T(x_i, \enc{\tau^p_i})$ holds.

  \emph{(C2) Hereditary completion.}  Assume
  $\T(w(z), \den{\hat I(\tup\alpha)}{w})$.  By \cref{def:MN} the second
  argument must lie in $\GTy \cup \{\ast\}$; the value
  $\hat I^{\MN}(w(\tup\alpha))$ lies in $\GTy$ only if
  $\mathrm{ok}_{\hat I}$ holds, i.e.\ all $w(\alpha_j) \in \GTy$ (a junk
  pair is neither a ground type nor $\ast$).  In that case
  $\T(w(\alpha_j), \ast)$ holds for each $j$; otherwise the assumption is
  false and the implication vacuous.
\end{proof}

\begin{theorem}[Soundness of the guarded translation]
  \label{thm:base-soundness-end}
  Let $\Env$ be an environment and $\phi_0$ a sentence.  If
  $\Ax \cup \Comp \cup \{\neg\enc{\phi_0}\}$ is unsatisfiable, then
  $\Th \entg \phi_0$.  In particular, if $\Env$ is consistent then
  $\Ax \cup \Comp$ is satisfiable.
\end{theorem}

\begin{proof}
  Contrapositive.  If $\Th \not\entg \phi_0$, pick a ground structure
  $\NN \models \Th$ with $\NN \not\entg \phi_0$.  By
  \cref{thm:base-soundness}, $\MN \models \Ax \cup \Comp$, and by
  \cref{lem:relativization}, $\MN \not\models \enc{\phi_0}$, i.e.\
  $\MN \models \neg\enc{\phi_0}$.  So the translated problem is
  satisfiable.  The second claim takes any $\NN \models \Th$ and observes
  $\MN \models \Ax \cup \Comp$.
\end{proof}

\begin{remark}
  \Cref{thm:base-soundness-end} quantifies over \emph{all} models of
  $\Th$, so its conclusion $\Th \entg \phi_0$ is classical semantic
  entailment from the source theory --- via \cref{prop:cic-bridge}, a
  consequence of the CIC environment under excluded middle.  This is the
  standard guarantee for hammers, whose final arbiter is proof
  reconstruction in the proof assistant; what the theorem rules out is the
  ATP refuting a translation of a goal that the source theory does not
  entail.  All optimization theorems below preserve exactly this
  guarantee.
\end{remark}
